\chapter[Discussion]{Discussion}\label{ch:sec:discussion}
The book transforms universal optimality from a probabilistic property
of deterministic encryption systems into an intrinsic geometric
property of weighted incidence structures.  The probabilistic
formulation (posterior uniformity on cloaks) motivates the theory, but
the weighted incidence structure becomes the intrinsic object, and the
Hilbert-space geometry is forced by its logarithmic linearization.
The development unfolds in five integrated layers.  The combinatorial
layer (Chapters~\ref{ch:sec:background}--\ref{ch:sec:graphs}) identifies weighted incidence structures as the
intrinsic invariant of universally optimal encryption, gives the exact
simple-cycle criterion for positive edge labels to admit a WIS
factorisation, establishes their gauge uniqueness, and completely
characterizes encoder-realizable
WIS by two intrinsic conditions: integrality of the induced
multiplicities and the linear balance equation $B\beta = K w$ on the
incidence graph.  The linearization layer (Chapter~\ref{ch:sec:hilbert-geometry}) takes logarithms
of the multiplicative WIS factorisation, revealing an additive
structure $\log n_{mc}=v_c-u_m$ that is the bridge to operator theory;
the canonical linear representation $\Phi(u,v)=Au-Cv$ is forced by this
representation, not chosen.  The analytic layer (Chapter~\ref{ch:sec:the-posterior-manifold}) then
unfolds inevitably from $\Phi$: the least-squares variational principle,
the normal operator (the bipartite incidence Laplacian), the posterior
Laplacian as the reduced Hessian, convexity of the reduced energy,
strict convexity modulo gauge, stability of the reconstruction, the
quotient geometry forced by posterior gauge symmetry, and the intrinsic
Euclidean distance from universal optimality.  The geometric layer pervades
the analytic one: the quotient space $V_I/\operatorname{Im}(C)$ is
determined by the posterior symmetry of Bayes' rule, and the quotient
norm provides the universal optimality defect.
A fifth, variational layer (Chapters~\ref{ch:sec:analytic-uod}--\ref{ch:sec:open}) develops the algebraic and
differential theory of the UOD, compares it with classical Fisher
geometry, establishes the Three-Level Theorem, and applies the
framework to eight concrete security and privacy problems.
One conceptual consequence deserves emphasis.  The realization problem
is not combinatorial once the multiplicity matrix is known: any
nonnegative integer matrix with constant row sums admits a deterministic
encoder by independently assigning ciphertext lists to each message
(Proposition~\ref{ch:prop:matrix-realization}).  The only genuinely
nontrivial obstruction therefore lies at the level of weighted incidence
structures, where integrality and the balance equation $B\beta=Kw$
completely characterize realizability.  This separates the combinatorial
encoding problem from the geometric structure of the weighted incidence
weights.
The book thus establishes two complementary results.  First, a
complete combinatorial classification of deterministic universally
optimal encryption systems.  Second, the discovery that every such
structure possesses a canonical Hilbert-space geometry whose analytic
constructions are developed in the subsequent parts.  They are
successive consequences of one canonical linearization.
This logical organization is reflected in the structure of the book:
Parts~I--II develop the core WIS theory, while Parts~III--V extend the
geometric analysis to information functionals, security measures, and
variational principles on the posterior manifold, unifying classical
entropy bounds and modern security estimates under a single geometric
framework.
\paragraph{Related work.}
The contributions of this book intersect four distinct research
traditions.  In each case we explain the existing literature and
identify what is genuinely new in the present work.
\subsection*{Perfect secrecy and posterior symmetry}\label{ch:sec:perfect-secrecy-and-posterior-symmetry}
The information-theoretic framework for secrecy was established by
Shannon~\cite{shannon1949}, who introduced the finite probabilistic
model of symmetric-key encryption used throughout the present work.
Shannon's notion of \emph{perfect secrecy} requires the posterior
distribution to equal the prior for every ciphertext:
\[
\Pr[M=m\mid C=c]=\Pr[M=m]
\qquad(\forall m\in\Mset,\; \forall c\in\Cset).
\]
This condition is well known to be equivalent to
$H(M)=H(M\mid C)$ and implies $|K|\ge|M|$.
The present work studies a different posterior symmetry:
universal optimality requires the posterior to be uniform over the
cloak of each ciphertext, i.e.~$\Pr[M=m\mid C=c]=1/|\cloak(c)|$
for every message in the cloak.  This condition does not compare the
posterior with the prior; it requires equiprobability of all messages
compatible with the observed ciphertext.  The two notions coincide
only when the prior is already uniform on every cloak.  Universal
optimality is therefore a structurally distinct symmetry property
within the same probabilistic framework, motivated by a different
structural objective.
Quantitative information-theoretic measures of security---equivocation
(as studied by Shannon~\cite{shannon1949} under the name of
``key equivocation''), mutual-information leakage, and min-entropy
leakage~\cite{smith2009foundations}---measure how much information
about the secret is revealed by the ciphertext.  The Universal
Optimality Defect introduced in this work serves a different purpose:
it measures the distance of a given encoder from the structural
condition of universal optimality, expressed as a Euclidean distance
in a canonically defined quotient Hilbert space.  The defect is
geometric rather than entropic, and its vanishing characterizes the
exact structural condition studied throughout this work.
\subsection*{Universal Optimality and Apollonian Cell Encoders}\label{ch:sec:universal-optimality-and-apollonian-cell-encoders}
The notion of universal optimality was introduced by Biondi,
Given-Wilson and Legay~\cite{biondi2018apollonian} in the context of
Apollonian Cell Encoders (ACE).  That paper constructs an explicit
family of deterministic encoders---the Apollonian cell encoders---and
proves that they achieve universal optimality with respect to every
entropic measure simultaneously.  The construction is algebraic and
relies on a specific representation of the encoder as a family of
cells with Apollonian packing structure.
The present work adopts the opposite viewpoint.  Rather than
constructing one family of encoders, it asks the structural question:
which deterministic encryption systems are universally optimal?  The
answer is a complete characterisation through weighted incidence
structures (Theorems~\ref{ch:thm:representation}
and~\ref{ch:thm:realizability}).  The logarithmic linearization, the
canonical mismatch operator $\Phi$, the Hilbert-space formulation, the
least-squares variational principle, the posterior Laplacian, the
quotient geometry, and the Universal Optimality Defect are all absent
from~\cite{biondi2018apollonian}.  Reference~\cite{biondi2018apollonian}
is an application paper that exhibits one family of universally optimal
encoders; the present work is the underlying mathematical theory that
characterizes the entire class.
\subsection*{Matrix scaling, biproportional fitting, and weighted incidence structures}\label{ch:sec:matrix-scaling-biproportional-fitting-and-weighted-inci}
The algebraic factorisation $n_{mc}=\beta(c)/w(m)$ strongly resembles
the multiplicative structure of classical matrix scaling (also known
as biproportional fitting, iterative proportional fitting, or the
RAS algorithm).  This literature begins with
Deming~and~Stephan~\cite{deming1940least}, who introduced iterative
proportional fitting for contingency tables, and continues with
Sinkhorn~\cite{sinkhorn1964relationship,sinkhorn1967concerning},
Bacharach~\cite{bacharach1970biproportional}, and the comprehensive
survey by Idel~\cite{idel2016review}.  These works study the existence
and uniqueness of positive diagonal scaling matrices $R$, $S$ such
that for a given nonnegative matrix $A$, the scaled matrix
$RAS$ satisfies prescribed row and column sums.
Although the multiplicative factorisation is algebraically similar,
the mathematical objective is fundamentally different.  Classical
matrix scaling starts from an existing matrix and searches for scaling
factors that satisfy prescribed marginal constraints.  The present
work, conversely, derives the multiplicative factorisation from Bayes'
rule applied to universally optimal deterministic encoders---it is a
consequence of the representation theorem, not an algorithmic
construction.  Weighted incidence structures are introduced as the
intrinsic mathematical objects, and the realizability theorem
(Theorem~\ref{ch:thm:realizability}) characterizes exactly which abstract
WIS arise from deterministic encoders through the two intrinsic
conditions of integrality and the balance equation $B\beta=Kw$.  The
subsequent operator-theoretic development (logarithmic linearization,
mismatch operator, variational principle, quotient geometry) has no
counterpart in the matrix balancing literature.
Uniqueness of the factorisation up to componentwise scaling
(Theorem~\ref{ch:thm:uniqueness}) is analogous to the well-known
uniqueness property of Sinkhorn scaling when the incidence graph is
connected~\cite{sinkhorn1964relationship}, although the proof here
follows from a short graph-theoretic argument (Lemma~\ref{ch:lem:edge-ratios}
and propagation along paths) rather than from the Perron--Frobenius
theorem.  The similarity is algebraic, not conceptual.
\subsection*{Graph Laplacians, graph potentials, and operator theory}\label{ch:sec:graph-laplacians-graph-potentials-and-operator-theory}
The logarithmic identity $\log n_{mc}=v(c)-u(m)$ places the problem
naturally in the theory of graph potentials on bipartite
graphs~\cite{bollobas1998modern,zaslavsky1998signed}.  Once the mismatch
operator $\Phi=[A\;-C]$ is identified, the subsequent analytical
constructions are standard consequences of Hilbert-space theory:
\begin{itemize}
\item the least-squares formulation and the normal equations
      $\Phi^{\mathsf T}\Phi\,x=\Phi^{\mathsf T}\omega$;
\item the observation that $\Phi^{\mathsf T}\Phi$ is the Laplacian
      of the bipartite incidence graph~\cite{chung1997spectral};
\item the elimination of the ciphertext potential $v$ as a Schur
      complement, which coincides with a
      Kron~\cite{kron1939tensor} reduction of the graph Laplacian;
\item convexity of the reduced energy, strict convexity modulo the
      gauge kernel $\ker\Phi$, and the Lipschitz stability estimate
      for the minimum-norm least-squares solution.
\end{itemize}
The claim of novelty is not these analytical tools themselves, each
of which is well known in its respective domain.  The novelty is the
proof that they arise \emph{canonically} from the logarithmic
representation of weighted incidence structures associated with
universally optimal encryption.  The mismatch operator is not an
externally imposed modelling choice; it is the unique linear operator
forced by the WIS representation (Theorem~\ref{ch:thm:mismatch}).  From
that operator, every subsequent construction---the normal operator,
the posterior Laplacian, convexity, strict convexity modulo gauge,
stability, the quotient geometry, and the Universal Optimality
Defect---follows through elementary Hilbert-space theory.  Thus the
operator-theoretic layer should be viewed as an intrinsic consequence
of the representation theorem rather than as an independently
postulated analytical framework.
Taken together, this work lies at the intersection of four research
traditions---information-theoretic cryptography, weighted incidence
combinatorics, matrix scaling theory, and graph operator theory---while
making a fundamentally different contribution from each: a complete
structural characterisation of deterministic universally optimal
encryption systems together with the canonical Hilbert-space geometry
and its associated Hilbert-space formalism.
\subsection*{Quantitative information flow}\label{ch:sec:quantitative-information-flow}
The security measures analysed in this book---min-entropy, Bayes
vulnerability, guessing entropy---are central to the field of
quantitative information flow (QIF), whose foundations were laid
by Smith~\cite{smith2009foundations}, Malacaria and colleagues
\cite{malacaria2007assessing, clark2007static}, and Palamidessi and
colleagues \cite{chatzikokolakis2005probability, alvim2012information,
alvim2023information}.  The WIS framework provides a unified geometric
treatment of these measures, showing that their piecewise-smooth
structure is a consequence of the same posterior ranking transitions
that govern all rank-based security functionals.
\paragraph{Final remark.}  This work's contribution is the complete
structural characterisation of deterministic universally optimal
encryption systems together with the canonical operator-theoretic
geometry forced by that characterisation.  Weighted incidence
structures are not merely descriptive invariants; together with two
intrinsic algebraic conditions they provide an equivalent mathematical
description of the entire class of universally optimal deterministic
encoders, and their logarithmic linearization induces a Hilbert-space
geometry whose analytic properties unfold inevitably.
Several questions remain open.  The most immediate are listed in
Chapter~\ref{ch:sec:open}.

\paragraph{What has been achieved.}
A complete classification of universally optimal deterministic encoders, a forced Hilbert-space geometry for every such encoder, universal Lipschitz and quadratic bounds for all smooth and piecewise-smooth security functionals, and a Three-Level Theorem that collapses the variational structure to at most three distinct values.

\paragraph{What becomes possible next.}
The open problems of Chapter~\ref{ch:sec:open} outline the natural extensions: randomised encoders, continuous message spaces, curvature, gradient flows, spectral theory, quantum generalisations.  Each of these builds on the geometric foundation laid here and extends the theory to new domains.

\endinput
